% Part: set-theory
% Chapter: spine
% Section: recursion

\documentclass[../../../include/open-logic-section]{subfiles}

\begin{document}

\olfileid{sth}{spine}{recursion}
\olsection{ট্রান্সফাইনাইট পুনরাবৃত্তির উপপাদ্যগুলি}

প্রথমে নিশ্চিত করতে হবে যে
\olref[valpha]{defValphas} সত্যিই একটি সফল সংজ্ঞা।
% BN-SRC-438: source's “offer a transfinite by (transfinite) recursion” is syntactically broken.
আরও সাধারণভাবে, ট্রান্সফাইনাইট পুনরাবৃত্তি দিয়ে
যেকোনো সংজ্ঞা দেওয়ার চেষ্টা সফল হয়—এটি প্রমাণ
করতে হবে। এই বিভাগের লক্ষ্য সেটিই।

\emph{সতর্কতা: বিষয়টি জটিল}। তবে সামগ্রিক শিক্ষা
সহজ: ট্রান্সফাইনাইট আবেশ এবং প্রতিস্থাপন মিলে
ট্রান্সফাইনাইট পুনরাবৃত্তির (বিভিন্ন রূপের) বৈধতা
নিশ্চিত করে।\footnote{মনে করিয়ে দিই: অন্য কথা
স্পষ্টভাবে না বললে সব সূত্র ও পদে পরামিতি থাকতে
পারে।}

\begin{defn}
ধরি, $\tau(x)$ একটি পদ, $f$ একটি ফাংশন এবং
$\alpha$ একটি ক্রমসংখ্যা। $f$-কে $\tau$-র জন্য
একটি \emph{$\alpha$-সন্নিকটায়ন} বলব যদি
$\dom{f} = \alpha$ এবং $(\forall \beta \in \alpha)f(\beta) = \tau(\funrestrictionto{f}{\beta})$
উভয়ই সত্য হয়।
\end{defn}

\begin{lem}[সীমাবদ্ধ পুনরাবৃত্তি]\ollabel{transrecursionfun}
যেকোনো পদ $\tau(x)$ ও ক্রমসংখ্যা $\alpha$-এর
জন্য $\tau$-র একটি অনন্য $\alpha$-সন্নিকটায়ন
আছে।
\end{lem}
\begin{proof}
দেখাব, যেকোনো $\gamma \leq \alpha$-এর জন্য
একটি অনন্য $\gamma$-সন্নিকটায়ন আছে।

প্রথমে অনন্যতা প্রতিষ্ঠা করি। ধরি, যথাক্রমে
$g$ ও $h$ হলো $\gamma$- এবং
$\delta$-সন্নিকটায়ন। তাদের যুক্তির উপর
ট্রান্সফাইনাইট আবেশে দেখা যায়,
যেকোনো $\beta \in \dom{g} \cap \dom{h} =
\gamma \cap \delta = \min(\gamma, \delta)$-এর
জন্য $g(\beta) = h(\beta)$। তাই সন্নিকটায়ন
বিদ্যমান থাকলে অনন্য এবং সাধারণ যুক্তিগুলিতে
তাদের মান একই।

অস্তিত্ব প্রতিষ্ঠার জন্য এবার
\olref[ordinals][opps]{simpletransrecursion}-এর
সরল ট্রান্সফাইনাইট আবেশ $\delta \leq \alpha$
ক্রমসংখ্যাগুলিতে প্রয়োগ করি।

শূন্য ফাংশনটি স্বভাবতই একটি
$\emptyset$-সন্নিকটায়ন।

$g$ যদি $\gamma$-সন্নিকটায়ন হয়, তবে $g \cup
\{\tuple{\gamma, \tau(g)}\}$ একটি
$\ordsucc{\gamma}$-সন্নিকটায়ন।

$\gamma$ সীমা ক্রমসংখ্যা হলে এবং প্রত্যেক
$\delta < \gamma$-এর জন্য $g_\delta$ একটি
$\delta$-সন্নিকটায়ন হলে, ধরি
$g = \bigcup_{\delta \in \gamma} g_\delta$।
বিভিন্ন $g_\delta$ সাধারণ যুক্তিতে একই মান দেয়
বলে এটি একটি ফাংশন। আর $\delta \in \gamma$
হলে $g(\delta) = g_{\ordsucc{\delta}}(\delta) =
\tau(\funrestrictionto{g_{\ordsucc{\delta}}}{\delta}) =
\tau(\funrestrictionto{g}{\delta})$।

এতে ট্রান্সফাইনাইট আবেশের প্রমাণ শেষ।
\end{proof}

ফাংশনের বদলে একটি \emph{পদ} সংজ্ঞায়িত করতে
দিলে আগের ফল থেকে $\alpha$ সীমাটি সরাতে
পারি। পরের ফল ও প্রমাণে $\sigma$ পদ হলে
$\funrestrictionto{\sigma}{\alpha} = \Setabs{\tuple{\beta,
\sigma(\beta)}}{\beta \in \alpha}$ নিই।

\begin{thm}[সাধারণ পুনরাবৃত্তি]\ollabel{transrecursionschema}
যেকোনো পদ $\tau(x)$-এর জন্য এমন একটি পদ
$\sigma(x)$ স্পষ্টভাবে সংজ্ঞায়িত করা যায় যে
যেকোনো ক্রমসংখ্যা~$\alpha$-এর জন্য
$\sigma(\alpha) = \tau(\funrestrictionto{\sigma}{\alpha})$।
\end{thm}

\begin{proof}
প্রত্যেক $\alpha$-এর জন্য
\olref{transrecursionfun} অনুযায়ী $\tau$-র
একটি অনন্য $\alpha$-সন্নিকটায়ন $f_\alpha$
আছে। $\sigma(\alpha)$-কে
$f_{\ordsucc{\alpha}}(\alpha)$ হিসেবে সংজ্ঞায়িত
করি। তাহলে:
	\begin{align*}
		\sigma(\alpha) &=
		f_{\ordsucc{\alpha}}(\alpha) \\&=
		\tau(\funrestrictionto{f_{\ordsucc{\alpha}}}{\alpha}) \\&=
		\tau(\Setabs{\tuple{\beta, f_{\ordsucc{\alpha}}(\beta)}}{\beta \in \alpha}) \\&=
		\tau(\Setabs{\tuple{\beta, f_{\ordsucc{\beta}}(\beta)}}{\beta \in \alpha})\\&=
		\tau(\funrestrictionto{\sigma}{\alpha})
	\end{align*}
শেষ সমতায় \olref{transrecursionfun}-এর মতো
প্রত্যেক $\beta < \alpha$-এর জন্য
$f_{\ordsucc{\beta}}(\beta) = f_{\ordsucc{\alpha}}(\beta)$
ব্যবহার করেছি।
\end{proof}
\noindent
লক্ষ করো, \olref{transrecursionschema} একটি
\emph{ছক}। $\sigma$ দিয়ে একটি ফাংশন—অর্থাৎ
বিশেষ একটি \emph{সেট}—সংজ্ঞায়িত হবে এমন
আশা করা যায় না। তা হলে $\dom{\sigma}$
সব ক্রমসংখ্যার সেট হতো, যা বুরালি-ফোর্তির
কূটাভাসের
(\olref[ordinals][basic]{buraliforti}) বিরোধী।

তবে এখনও দেখানো বাকি যে
\olref{transrecursionschema} আমাদের
$V_\alpha$-গুলির সংজ্ঞাকে বৈধতা দেয়। এটি
এখনই স্পষ্ট না হলেও ট্রান্সফাইনাইট
পুনরাবৃত্তির আরেকটি শেষ, সহজ রূপে স্পষ্ট হবে।

\begin{thm}[সরল পুনরাবৃত্তি]\ollabel{simplerecursionschema}
যেকোনো পদ $\tau(x)$ ও $\theta(x)$ এবং
যেকোনো সেট $A$-এর জন্য এমন একটি পদ
$\sigma(x)$ স্পষ্টভাবে সংজ্ঞায়িত করা যায় যে:
\begin{align*}
	\sigma(\emptyset) &= A\\
	\sigma(\ordsucc{\alpha}) &= \tau(\sigma(\alpha)) &&
		\text{যেকোনো ক্রমসংখ্যা }\alpha\text{-এর জন্য}\\
	\sigma(\alpha) &= \theta(\ran{\funrestrictionto{\sigma}{\alpha}})&&
	\text{যখন }\alpha\text{ সীমা ক্রমসংখ্যা}
\end{align*}
\end{thm}

\begin{proof}
প্রথমে নিচের মতো একটি পদ $\xi(x)$ সংজ্ঞায়িত
করি:
%t $\xi(x) = A$ যদি $x$ ক্রমসংখ্যা-সংজ্ঞাক্ষেত্রবিশিষ্ট ফাংশন না হয়;
%অন্যথায় নিই:
\[
	\xi(x) =
	\begin{cases}
			A &\text{যদি $x$ ফাংশন না হয় অথবা তার}\\
			&\hspace{1em}\text{সংজ্ঞাক্ষেত্র শূন্য বা অক্রমসংখ্যা হয়; অন্যথায়:}\\
			\tau(x(\alpha)) & \text{যদি $\dom{x} = \ordsucc{\alpha}$}\\
			\theta(\ran{x}) & \text{যদি $\dom{x}$ সীমা ক্রমসংখ্যা হয়}
	\end{cases}
\]
% BN-SRC-439: the source's default branch excludes the empty function although the proof needs xi(empty)=A.
\olref{transrecursionschema} অনুযায়ী এমন
একটি পদ $\sigma(x)$ আছে যে প্রত্যেক
ক্রমসংখ্যা $\alpha$-এর জন্য
$\sigma(\alpha) = \xi(\funrestrictionto{\sigma}{\alpha})$;
তদুপরি $\funrestrictionto{\sigma}{\alpha}$
হলো $\alpha$ সংজ্ঞাক্ষেত্রবিশিষ্ট একটি ফাংশন।
\olref[ordinals][opps]{simpletransrecursion}-এর
সরল ট্রান্সফাইনাইট আবেশে দেখাই, $\sigma$-র
চাওয়া ধর্মগুলি আছে।

প্রথমত, $\sigma(\emptyset) = \xi(\emptyset) = A$।

এরপর, $\sigma(\ordsucc{\alpha}) = \xi(\funrestrictionto{\sigma}{\ordsucc{\alpha}}) = \tau(\funrestrictionto{\sigma}{\ordsucc{\alpha}}(\alpha)) = \tau(\sigma(\alpha))$।

সবশেষে, $\alpha$ সীমা হলে
$\sigma(\alpha) = \xi(\funrestrictionto{\sigma}{\alpha}) = \theta(\ran{\funrestrictionto{\sigma}{\alpha}})$।
\end{proof}
\noindent
এখন \olref[valpha]{defValphas}-কে বৈধতা দিতে
শুধু $A
= \emptyset$, $\tau(x) = \Pow{x}$ এবং
$\theta(x) = \bigcup x$ নিলেই চলে। অবশেষে
$V_\alpha$-গুলির সংজ্ঞার ভিত্তি প্রতিষ্ঠিত হলো!{}

\end{document}
